\documentclass[12pt]{article}

% Assignment-manual style: Times text/math, restrained GitHub-like code blocks.
\usepackage[margin=1in]{geometry}
\usepackage{newtxtext,newtxmath}
\usepackage[T1]{fontenc}
\usepackage[utf8]{inputenc}
\usepackage{microtype}
\usepackage{amsmath,mathtools}
\usepackage{booktabs,longtable,array,tabularx}
\usepackage{xcolor}
\usepackage{listings}
\usepackage{enumitem}
\usepackage{hyperref}
\usepackage{fancyhdr}
\usepackage{lastpage}

\definecolor{githubink}{HTML}{24292F}
\definecolor{githubblue}{HTML}{0969DA}
\definecolor{githubgreen}{HTML}{1A7F37}
\definecolor{githubgray}{HTML}{57606A}
\definecolor{codebg}{HTML}{F6F8FA}
\definecolor{codeborder}{HTML}{D0D7DE}

\hypersetup{
  colorlinks=true,
  linkcolor=githubblue,
  citecolor=githubblue,
  urlcolor=githubblue,
  pdftitle={Assignment Handout: Competitive Data-Center Connection Planning},
  pdfauthor={Power Systems Course}
}

\lstdefinestyle{githubcode}{
  basicstyle=\ttfamily\small,
  backgroundcolor=\color{codebg},
  frame=single,
  rulecolor=\color{codeborder},
  framerule=0.5pt,
  columns=fullflexible,
  keepspaces=true,
  breaklines=true,
  showstringspaces=false,
  xleftmargin=0.5em,
  xrightmargin=0.5em
}
\lstset{style=githubcode}

\pagestyle{fancy}
\setlength{\headheight}{15pt}
\fancyhf{}
\lhead{Competitive Connection Planning}
\rhead{Assignment Handout}
\cfoot{\thepage\ of \pageref{LastPage}}

\newcommand{\code}[1]{\texttt{#1}}
\newcommand{\clioption}[1]{\lstinline|#1|}
\newcommand{\MW}{\,\mathrm{MW}}
\newcommand{\MWh}{\,\mathrm{MWh}}
\newcommand{\USD}{\$}

\title{\textbf{Assignment Handout}\\[0.4em]
  \large Competitive Data-Center Connection Planning Game}
\author{Course name}
\date{\today}

\begin{document}
\maketitle

\begin{center}
\fbox{\parbox{0.91\textwidth}{
\textbf{Assignment in one sentence.} You operate a data-center, choose
where and how much load to connect in a power grid, optionally build private
generation and storage, and maximize your score across three weekly matches.
}}
\end{center}

\tableofcontents
\newpage

\section{Assignment overview}

Each student group represents a data-center operator such as a cloud-computing
company. Your group submits a connection-and-resource plan. Submissions of
all groups will be collected and evaluated together using a DC optimal power
flow (\code{DCOPF}) dispatch model. \textbf{Groups are ranked by score.}

The classroom game has three weeks and six half-year submissions. Each week
has two submissions. Week 1 is practice: both of its submissions receive
feedback only and do not count. In Weeks 2 and 3, \textbf{both Submission 1 and
Submission 2 are scored and both count}. Your final ranking is the sum of four
scores: Week 2 Submission 1 + Week 2 Submission 2 + Week 3 Submission 1 +
Week 3 Submission 2.

\subsection{What you must do}
Your team must design and submit one complete half-year connection plan in the
released CSV format. The goal is to serve valuable computational load while
controlling energy, interconnection, generation, and storage costs.

\subsection{Deliverables}
\begin{enumerate}[leftmargin=*]
  \item one \code{UTF-8 CSV} file for your team;
  \item validator output showing that the canonical file passes the checker;
  \item a short strategy note explaining your main bus, load, generator, and
  storage choices and the trade-offs you considered;
  \item the exact final CSV and any code used to generate it.
\end{enumerate}
The instructor will announce the file name, submission location, deadline,
round assignment, and whether the strategy note is required for grading. Do
not submit a screenshot in place of the CSV.

\subsection{Assignment rules at a glance}
\begin{itemize}[leftmargin=*]
  \item Submit one canonical site-level CSV row for each selected site. The CSV has no
  time-period, load, or horizon rows.
  \item Use \code{bus IDs} from the supplied ACTIVSg2000 scenario block. Site
  capacities are specified in MW/MWh.
  \item The released calendar still determines the evaluator's half-year
  scenario, which uses six-hour periods (four per day): 2028 has 728/736 periods
  and 2027/2029 have 724/736 periods for January--June/July--December. Those
  periods are not CSV rows.
  \item \code{Generators} and \code{storage} are private to the submitting team.
  \item The active contract has no student-controlled \code{curtailment} decision.
  \item The instructor evaluates valid teams jointly and ranks teams by \code{score}.
\end{itemize}

\paragraph{Use of AI.} You are explicitly encouraged to use available AI tools while doing the
assignment, including ChatGPT, Codex, Claude Code, GitHub Copilot, and similar
tools. You remain responsible for checking generated code, understanding your
submission, respecting the CSV contract, and explaining your engineering
decisions. AI assistance is part of the intended computational workflow, not a
substitute for validation.

\subsection{Learning objectives}
\begin{itemize}[leftmargin=*]
  \item translate a planning idea into a machine-readable decision file;
  \item reason about spatial congestion, generator dispatch, and storage;
  \item understand why individually feasible plans can be jointly infeasible;
  \item use reproducible computational experiments and inspect diagnostics;
  \item communicate trade-offs among IT value, energy cost, and investment.
\end{itemize}

\section{Data and model: what is real and what is instructional}

The dispatch input is the public synthetic Texas ACTIVSg2000 model~\cite{activsg2000}, which
contains exactly \textbf{2,000 buses}, 3,206 branches, and 544 generator
records. The fixed topology is combined with a controlled three-year scenario
package built from hourly source records and released at a six-hour time step:
each game period is the mean of six consecutive source hours. The scenario
provides deterministic bus-level load; active
generator status, PMIN, and PMAX values come directly from the dataset. Because
ACTIVSg2000 does not provide technology labels or hourly availability data, the
active release does not apply synthetic wind/solar deratings. Each block is
validated period by period with DCOPF; the checked-in six-hour package
manifests record a passed validation for all six blocks.
It is an instructional scenario, not historical ERCOT telemetry or an official
ERCOT operating model.

The conceptual pipeline is:
\begin{center}
\fbox{\parbox{0.82\textwidth}{\centering
Public ACTIVSg2000 topology + controlled six-hour operating profiles\\
$\Downarrow$\\
Fixed ACTIVSg2000 buses, branches, generators, costs, and scenario demand\\
$\Downarrow$\\
DCOPF input and teaching scoring scenario\\
$\Downarrow$\\
Joint evaluation of student CSV files
}}
\end{center}

The following distinction is important:
\begin{itemize}[leftmargin=*]
  \item \textbf{Public data:} \code{bus}, \code{branch}, \code{generator}, and
  marginal-cost records from the ACTIVSg2000 source model.
  \item \textbf{Instructor preparation:} 
  \code{formatting}, \code{validation}, and the full-horizon economic profiles supplied with the round.
  \item \textbf{Instructional assumptions:} the student-owned investments,
  connection requests, \code{DCOPF} treatment, \code{IT value}, capital-cost coefficients,
  storage-cost coefficients, and the \code{fair continuous projection} policy.
\end{itemize}

The exercise should not be interpreted as an ERCOT planning recommendation or
as a forecast of actual interconnection feasibility.

\section{Operating horizon and decisions}

Each submission governs \textbf{one half-year with a six-hour time step}
($\Delta t=6$\,h, four periods per day). Each period is the mean of six
consecutive hourly source records:
\begin{center}
\begin{tabular}{lllll}
\toprule
Year & Round & Source hours & Six-hour periods & Zero-based period indices\\
\midrule
2028 & January--June & 4,368 & 728 & $0,\ldots,727$\\
2028 & July--December & 4,416 & 736 & $0,\ldots,735$\\
2027 or 2029 & January--June & 4,344 & 724 & $0,\ldots,723$\\
2027 or 2029 & July--December & 4,416 & 736 & $0,\ldots,735$\\
\bottomrule
\end{tabular}
\end{center}

The source year contains one full calendar year's hourly records, which the
released package averages into six-hour periods. A submission therefore has
access to roughly 2,000 possible buses and roughly 730 six-hour evaluation
states, but it does \emph{not} encode those states as CSV rows. The canonical
CSV lists only the sites used by your team; the evaluator combines those site
decisions with the selected six-hour scenario and economic artifacts. Every
energy quantity in the score (IT value, energy purchases, and generator
operating cost) is a power in MW multiplied by $\Delta t=6$\,h.

\subsection{Three-week, six-submission schedule}
\begin{center}
\small
\begin{tabularx}{0.98\textwidth}{l l l X}
\toprule
Week & Submission 1 & Submission 2 & Treatment\\
\midrule
1 & \code{2027\_jan\_jun} & \code{2027\_jul\_dec} & Practice week: feedback only; neither submission counts.\\
2 & \code{2028\_jan\_jun} & \code{2028\_jul\_dec} & Scored weekly match; Submission 1 and Submission 2 both count.\\
3 & \code{2029\_jan\_jun} & \code{2029\_jul\_dec} & Scored weekly match; Submission 1 and Submission 2 both count.\\
\bottomrule
\end{tabularx}
\end{center}

Within a week, the evaluator carries accepted sites and evaluator-assigned
generator and storage assets from Submission 1 into Submission 2. Storage
energy does not carry over: every storage device starts and ends each day
exactly half full, so each submission starts from the same state of charge.
Capital cost (site cost, generator capital, and
storage capital) is charged once, in the submission where the asset is first
built: Submission 1 pays for everything it builds, and Submission 2 pays only
for sites and assets newly added in Submission 2. Assets carried over from
Submission 1 are not charged again. Submission 1 has no energy-purchase cost;
Submission 2 pays for grid purchases at the joint-dispatch LMP. At the
beginning of the next week, the asset ledger resets; no asset carries between
weeks.

For each selected site $s$, submit a nonnegative connection capacity and
optional private generation and storage capacities. A team's resources serve
only that team. They cannot export energy to another team or become a shared
pool. The \code{bus\_id} must come from the supplied bus catalog.

\section{Mathematical model}

\subsection{Connection load}

Let $C_{i,s}\geq 0$ be the connection capacity submitted by team $i$ for site
$s$. The canonical CSV stores $C_{i,s}$ directly, together with optional
generator and storage capacities; it does not store a time-series load profile.
The evaluator uses the site-level plan when constructing each six-hour dispatch.

\subsection{Private generation}

Each positive generator capacity in the canonical CSV represents a team-built
generator. The seven-column CSV has no generator- or storage-ID columns; the
evaluator assigns deterministic private asset IDs internally from the team and
site IDs. These IDs are not lookup keys into the public generator catalog. The
evaluator applies the
same convex piecewise-linear operating-cost curve to every team-built generator
at every site: the first 25\% of capacity costs \$80/MWh, the next 50\% costs
\$150/MWh, and the final 25\% costs \$250/MWh. You submit the capacity and the
evaluator chooses the dispatch in every six-hour period. For a generator with capacity $G_{i,g}$,
dispatch satisfies
\[
  0\le g_{i,g,t}\le G_{i,g}.
\]
The generator is physically attached to the submitted bus and is private to your
team. It is not an existing public generator and does not inherit a public
generator ID, PMIN, PMAX, or marginal cost.

\subsection{Private storage}

For storage device $b$, submit power capacity $P_{i,b}$ in MW and energy
capacity $B_{i,b}$ in MWh. The evaluator assigns the device's deterministic
private asset ID internally; the student CSV contains no \code{storage\_id}
field. Because efficiency is ideal and storage has no operating cost, the
evaluator represents each device by one continuous signed net-power variable
$p_{i,b,t}$ per six-hour period (positive discharges, negative charges), so the
dispatch remains a linear program with no binary charge/discharge mode
variable. With $\Delta t=6$\,h, the stored energy evolves as
\[
  e_{i,b,t}=e_{i,b,t-1}-\Delta t\,p_{i,b,t},
\]
with
\[
  0\le e_{i,b,t}\le B_{i,b},\qquad
  -P_{i,b}\le p_{i,b,t}\le P_{i,b},\qquad
  \Delta t\,p_{i,b,t}\le e_{i,b,t-1}.
\]
Storage operates on an exact daily cycle: each day consists of four
consecutive six-hour periods, every day starts with $e=0.5B_{i,b}$, and every
day must end with exactly $e=0.5B_{i,b}$. Stored energy never carries from one
day to the next, nor from one submission to another. Within a day, energy
charged in one period can be discharged in a later period of the same day.
Students submit storage capacities only. The grid operator chooses charging and
discharging in the centralized dispatch. Storage is private to the submitting
team: charging appears as additional team demand and discharging reduces that
team's grid purchase; it cannot serve any other team's load. The evaluator
assigns the storage asset ID from the team/site identity.

\paragraph{Storage cost calibration.} The active half-year scoring files use
the following instructional coefficients:
\begin{align*}
 C^{\mathrm{storage,cap}}_{i,b}={}&
 (\USD400{,}000/\MW)P_{i,b}
 +(\USD250{,}000/\MWh)B_{i,b}.
\end{align*}
For a four-hour device ($B_{i,b}=4P_{i,b}$), the two capital terms correspond to
\USD1.4m/MW of installed power capacity. Storage operating cost is not charged.
These are rounded teaching coefficients informed by U.S. battery-storage cost
benchmarks, not a project quotation or an ERCOT tariff~\cite{eia2024bess,
nrel2023bess}. Storage capital is charged once, in the submission where the
device is first built.

\subsection{DCOPF}

For each six-hour period, the instructor solves a linear \code{DCOPF} on the prepared Texas
network. Network balance, generator limits, branch thermal limits, fixed base
demand, and the centralized team-dedicated storage model are enforced. Load is
inelastic in the ordinary feasible case: there is no integer unit commitment
and no ordinary load-shedding variable. The daily storage cycle couples the
four periods of each day; the solver, not the student, selects the storage
schedule.

\section{CSV submission format}\label{sec:csv}

Submit one \code{UTF-8 CSV} file per team. The canonical header is exactly:

\begin{lstlisting}
team_id,site_id,bus_id,connection_capacity_mw,generator_capacity_mw,storage_power_mw,storage_energy_mwh
\end{lstlisting}

Each data row is one site. The file has no \code{period}, \code{load\_mw},
\code{submission\_id}, \code{generator\_id}, or \code{storage\_id} columns and
contains no horizon rows. The \code{site\_id} is a team-defined site
label; \code{bus\_id} must come from the supplied scenario block.

Example:
\begin{lstlisting}
team_id,site_id,bus_id,connection_capacity_mw,generator_capacity_mw,storage_power_mw,storage_energy_mwh
team-blue,site-north,1001,40,20,10,40
team-blue,site-south,1002,15,0,0,0
\end{lstlisting}

Values use MW, MWh, and decimal notation. The \code{validator} checks the exact
seven-column header, site identifiers, scenario bus identifiers, finite
nonnegative capacities, and the storage ratio.

\textbf{Before submitting, run the provided \code{CSV validator}.} From the project
root, run with the released ACTIVSg2000 scenario-block path:
\begin{lstlisting}
PYTHONPATH=src python scripts/check_submission.py \\
  --file team-blue.csv
\end{lstlisting}
The \code{checker} is also available after installation as
\code{ercot-game-check-csv}. Add \clioption{--artifact-dir} when you want the checker
to verify scenario bus IDs; \clioption{--year}, \clioption{--round}, and \clioption{--horizon}
are optional historical/report context. The validator is a convenience check;
the instructor evaluator remains authoritative.

\textbf{Historical compatibility note.} The retired ten-column hourly format
with \code{period} and \code{load\_mw} is retained only for explicitly marked
legacy fixtures and compatibility tests. Do not use it for a current submission.
The legacy format's asset-ID fields are historical compatibility fields, not part
of the active student contract.

\subsection{Complete small example}
The following is a complete canonical two-site example. It is intentionally
small for format testing; a course submission may contain more site rows but
still never contains horizon rows:
\begin{lstlisting}
team_id,site_id,bus_id,connection_capacity_mw,generator_capacity_mw,storage_power_mw,storage_energy_mwh
team-demo,site-north,1001,4,2,1,4
team-demo,site-south,1002,2,0,0,0
\end{lstlisting}
The second site has no generation or storage. The half-year six-hour scenario is
selected by the instructor and is not expanded into repeated rows in this file.

\subsection{Checker troubleshooting}
\begin{itemize}[leftmargin=*]
\item \textbf{Unknown bus:} copy the \code{bus\_id} from the supplied ACTIVSg2000
 bus catalog.
  \item \textbf{Duplicate site:} give every site row a distinct \code{site\_id}.
  \item \textbf{Storage ratio:} use storage energy equal to four times storage
  power, or set both storage fields to zero.
  \item \textbf{Malformed CSV:} save as UTF-8 CSV and inspect commas,
  quotation marks, and the exact header.
\end{itemize}

\section{How joint evaluation works}

The instructor first rejects \code{malformed CSV} files in isolation; a malformed file
does not prevent other teams from being evaluated. Review the exact file
contract in Section~\ref{sec:csv} before diagnosing a rejection.

Valid plans are then evaluated jointly. If the joint plan is \code{feasible}, each
team is scored using its cleared load in every six-hour period, private resource dispatch, and
the announced scoring vectors.

If the joint plan is \code{infeasible}, the evaluator applies the \code{fair
continuous projection}, discussed in Section~\ref{sec:projection}. No first-come-first-serve ordering is used. The projected
load is then used for the score. The capital costs associated with submitted
sites, generators, and storage remain part of the score. See Section~\ref{sec:score} for the score calculation.

\section{Fair continuous projection}\label{sec:projection}

If the valid teams' full request is jointly infeasible, the evaluator applies
fixed anti-manipulation caps before searching for a projected dispatch. Each
site is capped at $100\MW$, each team at $500\MW$, and in a multi-team case the
largest team request is capped at the smallest positive team request. These
caps are instructor-side policy; students cannot tune them.

Let $d_{i,s,t}$ be the capped requested load of team $i$. The evaluator chooses
a separate service fraction $\alpha_i\in[0,1]$ for each team and evaluates
$\widehat d_{i,s,t}=\alpha_i d_{i,s,t}$. It first maximizes total served load.
It then restricts the search to solutions within
$\max(10^{-6},10^{-6}L^*)$ MWh of that optimum $L^*$ and maximizes
\[
  \sum_{i:d_i>0}\log(\alpha_i d_i),
\]
which is proportional-fair tie-breaking. Positive valid requests have a
minimum service fraction of $10^{-6}$. The result records each team's fraction,
requested load, projected load, and the cap policy.

This policy keeps valid teams eligible, prevents an oversized request from
dominating the projection, and produces a physically feasible joint dispatch.
It is a teaching rule, not an ERCOT operating procedure~\cite{kelly1998charging,
jacubowicz2022fairness,gebbran2021fairder}.

\section{Scoring}\label{sec:score}

For team $i$, the \code{score} is
\[
\begin{aligned}
S_i={}V_i^{\mathrm{IT}}-C_i^{\mathrm{energy}}
-C_i^{\mathrm{gen,op}}-C_i^{\mathrm{gen,cap}}-C_i^{\mathrm{storage,cap}}
-C_i^{\mathrm{site}}.
\end{aligned}
\]

The components are:
\begin{itemize}[leftmargin=*]
  \item $V_i^{\mathrm{IT}}$: cleared load multiplied by the explicit full-horizon
  \code{IT value} vector and by $\Delta t=6$\,h. The active AI-intensive
  profile uses \USD6500/MWh in every six-hour period;
  \item $C_i^{\mathrm{energy}}$: for Submission 2, grid purchase (MW) multiplied by
  the nonnegative LMP from the selected joint dispatch at the selected bus and
  six-hour period and by $\Delta t=6$\,h. Submission 1 has no energy-purchase cost;
  \item $C_i^{\mathrm{gen,op}}$: team-generator dispatch (MW) times
  $\Delta t=6$\,h, charged at the common 25\%/50\%/25\% curve of
  \USD80/150/250 per MWh;
  \item $C_i^{\mathrm{gen,cap}}$: generator MW newly built in this submission
  multiplied by \USD1.5m/MW;
  \item $C_i^{\mathrm{storage,cap}}$: storage newly built in this submission
  charged at \USD400,000/MW plus \USD250,000/MWh, with no storage operating cost;
  \item $C_i^{\mathrm{site}}$: connection capacity newly built in this
  submission multiplied by \USD25m/MW.
\end{itemize}
Capital terms are charged once, in the submission where a site or asset is
first built. In Submission 1 they cover everything you build; in Submission 2
they cover only sites and assets you add in Submission 2, because Submission 1
assets carry over without a second charge. Everything resets between weeks.
An invalid submission scores 0.

Your final ranking is the sum of the four scored submissions:
\[
  S_i^{\mathrm{final}}
  = S_i^{\mathrm{W2,S1}} + S_i^{\mathrm{W2,S2}}
  + S_i^{\mathrm{W3,S1}} + S_i^{\mathrm{W3,S2}}.
\]
Week 1 scores are practice feedback and are not included.

The active instructional coefficients are listed below. The selected rule file
and economic-profile artifacts remain authoritative for a released round.

\begin{center}
\small
\begin{tabularx}{0.98\textwidth}{>{\ttfamily\raggedright\arraybackslash}p{0.24\textwidth} >{\raggedright\arraybackslash}p{0.27\textwidth} X}
\toprule
Parameter & Active value & Basis\\
\midrule
Generator capital & \USD1.5m/MW & Charged once, in the submission where the generator is first built; rounded near EIA's 2024 internal-combustion-engine cost~\cite{eia2024generatorcosts}.\\
Site cost & \USD25m/MW & AI-data-center facility-capital proxy, not a universal interconnection tariff; informed by JLL and Morgan Stanley~\cite{jll2026dc,morganstanley2024dc}.\\
Generator operating cost & Piecewise \$80/150/250 per MWh at 25/50/25\% & Source-informed colocated natural-gas/diesel-standby curve; exact mapping is instructional~\cite{eia2026lcoe,eia2025heatrate}.\\
Energy price & Submission 2 joint-dispatch LMP & Locational price at the selected bus and six-hour period; Submission 1 has no energy-purchase cost.\\
IT value & \USD6500/MWh in every period & Full-horizon AI-intensive opportunity value, rounded from recent compute-value evidence; not a direct cloud tariff conversion~\cite{computeheat2026,iea2026ai}.\\
Storage power capital & \USD400,000/MW & Rounded 2028 moderate projection component~\cite{nrel2023bess}.\\
Storage energy capital & \USD250,000/MWh & Rounded 2028 moderate projection component; together with power capital gives \USD1.4m/MW for a four-hour system~\cite{nrel2023bess}.\\
Storage operation & Not charged & The active evaluator has no storage operating-cost term.\\
Invalid file score & 0 & Deterministic grading policy, not a market quantity.\\
\bottomrule
\end{tabularx}
\end{center}

\subsection{How to reason about your score}
The evaluator first determines how much of your requested load is served. It
then adds IT value for served load and subtracts energy, site, generator, and
storage costs. A larger plan is therefore not automatically better: an
additional site may increase IT value while also increasing energy,
congestion, and capital costs. A generator or storage device can reduce grid
purchases but still lower the final score if its costs exceed the avoided
energy cost.

Use the full-horizon \code{IT-value} vector and the locational \code{price}
artifact supplied with the round when estimating a score. The current IT-value
vector is constant by policy; the energy price is not a student-chosen value.

\section{Practical workflow}
\begin{enumerate}[leftmargin=*]
  \item Read the weekly schedule and identify the evaluator's scenario block.
  \item Inspect the ACTIVSg2000 bus and generator tables supplied by the instructor.
  \item Choose one or more buses and specify site capacities in the canonical CSV.
  \item Decide whether private generation or storage improves your score.
  \item Generate the \code{CSV} and run the \code{validator}.
  \item Keep the exact submitted file and validator output for reproducibility.
  \item Submit both weekly files before the announced deadlines. In Weeks 2
  and 3 both submissions count, so treat Submission 1 as a scored plan too;
  Submission 2 may add new sites and assets on top of it.
\end{enumerate}

\subsection{What to inspect after evaluation}
Read the result in this order:
\begin{enumerate}[leftmargin=*]
  \item feasibility or projected service status and your team-specific service fraction;
  \item total score and rank;
  \item cleared load, grid purchases, and resource costs;
  \item binding or limiting network diagnostics;
  \item storage state and generator dispatch;
  \item raw audit output only when debugging.
\end{enumerate}

\subsection{Final submission checklist}
Before uploading, confirm:
\begin{itemize}[leftmargin=*]
  \item the file is \code{UTF-8 CSV} with the exact header;
  \item each data row has a distinct site ID and the exact seven-column header;
  \item all bus IDs come from the supplied catalog;
  \item all capacities are finite and nonnegative, with storage energy equal to
    four times storage power when storage is present;
  \item the \code{checker} passes without requiring horizon context arguments;
  \item you retained the exact submitted file and checker output;
  \item your strategy note explains the main economic and physical trade-offs.
\end{itemize}

\section{Troubleshooting and frequently asked questions}
\begin{description}[style=nextline,leftmargin=0pt]
  \item[How are generator and storage assets identified?]
  The active seven-column CSV has no asset-ID columns, including no generator-ID
  or storage-ID columns. Submit capacities only; the evaluator assigns
  deterministic private asset IDs from the team and site IDs. You cannot choose
  the operating-cost curve: every team-built
  generator uses the common piecewise curve specified in Section~\ref{sec:score},
  and the evaluator dispatches it.

  \item[Can my generator or storage serve another team?]
  No. The grid operator may schedule your storage, but storage discharge and
  team-generator output can reduce only your team's load. They cannot be shared
  with another team.

  \item[What happens if the joint plan is infeasible?]
  The evaluator applies the capped fair projection in Section~\ref{sec:projection},
  using a separate service fraction for each valid team. The projection first
  maximizes total service and then uses proportional-fair tie-breaking.

  \item[What carries from Submission 1 to Submission 2?]
  Accepted sites and evaluator-assigned generators and storage assets persist
  within the week. Stored energy does not: every storage device starts and ends
  each day exactly half full. Capital is charged once, in the
  submission where an asset is first built: carried-over Submission 1 assets
  are not charged again in Submission 2, and Submission 2 pays capital only
  for newly added sites and assets. The ledger resets at the start of
  the next week.

  \item[Can I schedule storage myself?]
  No. Submit only storage power and energy capacities. The centralized grid
  optimizer schedules charge, discharge, and SOC subject to the four-hour,
  team-dedicated constraints and the exact daily cycle (half full at the start
  and end of every day).

  \item[Why is my score negative?]
  Served-load value may be smaller than energy, site, generator, or storage
  costs. Inspect the detailed score components before changing the plan.

  \item[Which result is official?]
  The instructor-generated result and audit files determine the official
  score. Week 1 is feedback-only; the final ranking is the sum of Submission 1
  and Submission 2 scores from Weeks 2 and 3 (four scores in total). Local calculations are for debugging and
  strategy development.

  \item[Is the scenario package validated?]
  Yes. The manifests of the released six-hour package
  (\code{activsg2000\_scenarios\_v2\_6h}) record a passed DCOPF validation for
  all six half-year blocks. Running the evaluator or the student preview
  locally still requires the \code{highspy} solver and the prepared artifacts.
  The official score always uses the full six-hour half-year horizon; the
  student preview samples only every sixth period (every 36 hours) and leaves
  storage idle, so it is a screening estimate, not an official score.
\end{description}

\section{References}
\bibliographystyle{plain}
\bibliography{references}

\end{document}
